\section{Conclusion}

Maintaining a connected laboratory requires collaborative work across interfaces, software versions, physical resources, material state, scientific constraints, and organizational boundaries. Public technical discussions make part of that work visible. Participants articulate the system omitted from an initial question, coordinate repair across software history and irreversible physical action, navigate porous boundaries between public and private support, and sometimes convert local diagnosis into durable community infrastructure.

The boundary-centered account developed here connects articulation work, infrastructure, and repair. Breakdowns expose dependencies, although visibility alone does not make knowledge durable. Durability requires provenance, applicability, validation, ownership, and a correction path. Recovery requires a record of uncertainty and final material disposition. Public support requires return paths when resolution moves into restricted channels.

These findings suggest that technical communities and automation platforms should be designed together. Progressive question interfaces, linked run-event records, bounded public resolution records, evidence grades, and maintenance-aware knowledge indexes can support cooperative diagnosis and repair. Their value should be evaluated through participation, clarification burden, safe-state time, public return, recurrence, staleness, and conversion into maintained artifacts. The connected laboratory is sustained not only by interoperable machines, but by infrastructures through which people make its dependencies inspectable, repairable, and correctable.
